\scanpage{54}
\mspage{IV}
% archive foliation: 53
(i) Pour tout $X$ de t.f. sur $S$, $H^1(X,G)$ est fini.

(ii) Pour tt sous-schéma \deleted{\unsure{affine}} normal
\inserted{\unsure{intègre affine}}
\deleted{\unsure{de }$S$}
$S_0$ de $S$, \deleted{\unsure{et}} tt \unsure{rev.} fini étale
\inserted{affine}
$S'$ de $S_0$, tout schéma lisse de t. fini $X$ sur $S'$, à fibres géom. connexes non vides, $H^1(\unsure{X},G)$ est fini.

(iii) Pour tt $S',X$ comme dessus, $S'\neq\varnothing$, il existe un ouvert non vide $U$ de $S'$ tel que $H^1(X|\unsure{U'|U},G)$ \unsure{soit} fini.

Il \uuline{suffit} pour ceci qu'on ait ceci :

(iv) a) Pour tout $s\in S$, et tt schéma $\bar X$ de type fini lisse affine connexe sur $\overline{k(s)}$,
\[
H^1(\bar X,G)\text{ est fini}
\]

b) Pour tout \deleted{\unsure{$S$}} \deleted{\unsure{un}} schéma affine normal intègre $S'$ \deleted{de $S$},
\inserted{comme dans (ii)}
et tt sous groupe $H$ de $G\times G$,
\deleted{\illegible{1}}
$H^1(S',H)$ est fini.

Lorsque $G$ est commutatif, on peut affaiblir b) en y remplaçant $H$ par $G$.

\figure{\begin{tikzpicture}
\draw (0,0)--(3.2,0);
\end{tikzpicture}}

\uline{Dém.} Equivalence de (i) à (iii) résulte aussitôt de lemme 3. Pour prouver qu'\unsure{ils} sont entrainés par (iv), on doit utiliser en plus lemme 1.
